\section{Conclusion}
Ce stage a permis d'explorer et de mettre en pratique des compétences variées, alliant pédagogie, conception d’outils éducatifs et compétences techniques. L'objectif initial — \textit{Apprendre à discuter avec l'IA} — a été atteint grâce à une approche structurée en deux parties : une dimension théorique (compréhension des modèles de langage via des énigmes) et une dimension pratique (interaction avec un modèle IA, révélant ses limites et biais).

Les défis rencontrés, comme l'équilibre entre complexité et accessibilité des énigmes ou la gestion des biais du modèle utilisé, ont renforcé notre compréhension des enjeux liés à la vulgarisation scientifique. L'animation de l'atelier auprès des élèves a confirmé leur curiosité et leur réceptivité, tout en soulignant un besoin urgent d'intégrer une sensibilisation à l'IA dans les programmes scolaires.

Ce projet ouvre des perspectives pour améliorer :
\begin{itemize}
    \item Les supports pédagogiques (ex : outils interactifs plus intuitifs).
    \item La formation des enseignants sur les limites et usages éthiques de l'IA.
\end{itemize}

Cette expérience a été formatrice, tant sur le plan technique que relationnel, et renforce mon intérêt pour la médiation scientifique et l'éducation aux technologies.
